% ============================================================
%  5.4  Desarrollo del módulo de captura y registro de intentos
%  Pegar a continuación de la sección 5.3 en main.tex
% ============================================================

\subsection{Desarrollo del módulo de captura y registro de intentos}

En esta sección se describe la implementación del módulo de captura,
cuya responsabilidad es obtener una imagen de la escritura del alumno
y transmitirla al servidor para su posterior análisis. El módulo se
concretó en dos componentes de interfaz independientes, cada uno
diseñado para un modo de entrada diferente:
\textit{CapturaEscritura}, que permite al alumno fotografiar un
trabajo realizado en papel, y \textit{LienzoDigital}, que ofrece un
lienzo interactivo para escribir directamente sobre la pantalla.
Ambos componentes comparten el mismo contrato de comunicación con el
\textit{backend}: una petición \texttt{POST} al endpoint
\texttt{/api/intentos/registrar} con la imagen codificada en
Base64 y el identificador del ejercicio asignado.

% ------------------------------------------------------------
\subsubsection{Implementación del componente \textit{CapturaEscritura}}
% ------------------------------------------------------------

El componente \textit{CapturaEscritura} atiende el caso en que el
alumno realiza el ejercicio de forma manuscrita en papel y necesita
digitalizar su trabajo mediante la cámara o la galería del
dispositivo. Su implementación resuelve tres retos prácticos: la
heterogeneidad de dispositivos (móvil versus escritorio), la
variabilidad de calidad de las imágenes obtenidas, y la necesidad de
ofrecer herramientas de edición básica antes del envío.

\subsubsection{Detección del dispositivo y apertura de la cámara}

La lógica de captura varía según el tipo de dispositivo. En
dispositivos móviles, el sistema operativo expone una interfaz nativa
de cámara a través del atributo \texttt{capture="environment"} del
elemento \texttt{<input type="file">}, la cual ofrece mayor
compatibilidad que las APIs web. En equipos de escritorio, en cambio,
se utiliza la API \texttt{MediaDevices.getUserMedia} para abrir una
vista de la cámara directamente dentro de la interfaz de la
aplicación. La función \texttt{handleTomarFotoClick} encapsula esta
bifurcación detectando el agente de usuario mediante una expresión
regular; si el resultado es positivo activa el \texttt{<input>} nativo,
de lo contrario solicita el stream de video con la cámara trasera como
preferencia y una resolución ideal de 1\,280 píxeles de ancho:

\begin{lstlisting}[language=Python, caption={Detección de dispositivo y apertura de cámara}, label={lst:deteccion_dispositivo}]
const handleTomarFotoClick = () => {
  const isMobile =
    /iPhone|iPad|iPod|Android/i.test(navigator.userAgent);
  if (isMobile) cameraInputRef.current.click();
  else iniciarCamaraWeb();
};

const iniciarCamaraWeb = async () => {
  if (!navigator.mediaDevices?.getUserMedia) {
    setError(
      'La camara no esta disponible. Usa "Subir Archivo".'
    );
    return;
  }
  try {
    const mediaStream = await navigator.mediaDevices
      .getUserMedia({
        video: { facingMode: 'environment',
                 width: { ideal: 1280 } }
      });
    setStream(mediaStream);
    setMostrarCamaraWeb(true);
  } catch {
    setError('No se pudo acceder a la camara.');
  }
};
\end{lstlisting}

Cuando la vista de cámara está activa, se muestra sobre ella una
plantilla guía (\textit{PlantillaGuia}) que delimita el área de
encuadre recomendada. Al presionar el botón de captura, el frame
actual del elemento \texttt{<video>} se transfiere a un canvas
temporal mediante \texttt{drawImage} y se convierte a \textit{data URL}
en formato JPEG. Una vez obtenida la imagen, el stream se libera
deteniendo todas sus pistas de video, evitando así que el indicador
de cámara del dispositivo permanezca activo.

\subsubsection{Preprocesamiento y normalización de la imagen}

Las imágenes capturadas o seleccionadas desde la galería pueden tener
dimensiones y pesos muy variados. Para garantizar que el payload
enviado al servidor sea manejable sin degradar la información útil
para el OCR, la función \texttt{procesarArchivo} valida primero el
tipo MIME y rechaza archivos que superen los 10\,MB, con lo que
proporciona retroalimentación inmediata antes de iniciar cualquier
transferencia. Si la imagen es válida, se aplica un redimensionamiento
proporcional que limita el lado mayor a 1\,200 píxeles, y el resultado
se recodifica como JPEG con calidad 0.8:

\begin{lstlisting}[language=Python, caption={Redimensionamiento y normalización de la imagen capturada}, label={lst:redimensionamiento}]
const procesarArchivo = (file) => {
  if (!file.type.startsWith('image/')) {
    setError('Solo se admiten imagenes.'); return;
  }
  if (file.size > 10 * 1024 * 1024) {
    setError('La imagen es demasiado grande. Maximo 10MB.');
    return;
  }
  const reader = new FileReader();
  reader.onloadend = () => {
    const img = new Image();
    img.onload = () => {
      const canvas = document.createElement('canvas');
      let { width, height } = img;
      const MAX = 1200;
      if (width > height && width > MAX) {
        height *= MAX / width; width = MAX;
      } else if (height > MAX) {
        width *= MAX / height; height = MAX;
      }
      canvas.width = width; canvas.height = height;
      canvas.getContext('2d')
            .drawImage(img, 0, 0, width, height);
      setImagenPreview(canvas.toDataURL('image/jpeg', 0.8));
    };
    img.src = reader.result;
  };
  reader.readAsDataURL(file);
};
\end{lstlisting}

\subsubsection{Herramientas de edición: rotación y recorte}

Una imagen capturada con el dispositivo en orientación incorrecta o
con exceso de contexto visual alrededor de la escritura puede afectar
negativamente el reconocimiento OCR. Para mitigar este riesgo, el
componente expone dos herramientas de edición que operan íntegramente
en el cliente.

\textbf{Rotación.} La función \texttt{rotarImagen} aplica una rotación
de 90° en sentido horario intercambiando las dimensiones del canvas y
aplicando una transformación afín sobre el contexto 2D mediante
\texttt{ctx.translate} y \texttt{ctx.rotate(Math.PI / 2)}.
El resultado se serializa nuevamente a JPEG y sustituye la vista
previa actual.

\textbf{Recorte interactivo.} Al activar el modo de recorte, se
superpone sobre la imagen una región de selección con ocho manejadores
de arrastre (cuatro esquinas y cuatro puntos medios). La lógica de
arrastre escucha los eventos \texttt{pointermove} y \texttt{pointerup}
sobre el objeto \texttt{window}, de modo que el desplazamiento
continúa aunque el cursor salga momentáneamente del manejador. El
estado de la selección se almacena en coordenadas relativas al
elemento \texttt{<img>} renderizado; al confirmar el recorte, dichas
coordenadas se escalan a las dimensiones naturales de la imagen
mediante los factores \texttt{naturalWidth\,/\,clientWidth} y
\texttt{naturalHeight\,/\,clientHeight}, garantizando precisión
independientemente del tamaño de visualización.

\subsubsection{Envío del intento al servidor}

Una vez satisfecha con la imagen, el alumno la envía pulsando el botón
de confirmación. El componente recupera el token JWT almacenado en
\texttt{localStorage}, construye el cuerpo de la petición con el
identificador del ejercicio y la imagen en formato \textit{data URL},
y realiza la solicitud \texttt{POST} al endpoint de registro. Ante
una respuesta exitosa se limpia el estado del formulario y se invoca
el \textit{callback} \texttt{alTerminar}, que notifica al componente
padre para que muestre la confirmación o navegue a la pantalla de
resultados.

% ------------------------------------------------------------
\subsubsection{Implementación del componente \textit{LienzoDigital}}
% ------------------------------------------------------------

El componente \textit{LienzoDigital} ofrece un lienzo de dibujo basado
en el elemento \texttt{<canvas>} de HTML5 para que el alumno escriba
directamente sobre la pantalla con ratón o con el dedo en dispositivos
táctiles. Incluye herramientas de lápiz y goma con grosor ajustable,
un mecanismo de deshacer por historial de estados, y un modo de
pantalla completa que amplía el área de trabajo.

\subsubsection{Inicialización y ciclo de dibujo}

El canvas se dimensiona al montarse el componente para que ocupe el
área disponible de su elemento padre; esta inicialización se repite
cada vez que el modo de pantalla completa cambia, garantizando que
el área de dibujo llene el contenedor sin deformaciones. Los atributos
\texttt{lineCap} y \texttt{lineJoin} se fijan en \texttt{'round'} para
suavizar extremos y uniones de cada trazo, produciendo una apariencia
más cercana a la escritura real.

Para garantizar compatibilidad táctil, la función
\texttt{obtenerCoordenadas} extrae la posición del puntero desde
\texttt{e.touches[0]} en eventos táctiles o desde
\texttt{e.clientX} / \texttt{e.clientY} en eventos de ratón. Las
coordenadas resultantes se escalan con el factor
\texttt{canvas.width\,/\,rect.width}, que compensa la diferencia entre
el tamaño lógico del canvas y el tamaño visual renderizado por CSS,
evitando desplazamientos del trazo respecto al punto de contacto.

La diferenciación entre lápiz y goma se realiza a través de la
propiedad \texttt{globalCompositeOperation} del contexto 2D. Con el
lápiz activo se usa \texttt{'source-over'}, que superpone nuevos
píxeles sobre los existentes; con la goma se usa
\texttt{'destination-out'}, que elimina los píxeles del canvas de
forma no destructiva sin reemplazarlos por ningún color de fondo:

\begin{lstlisting}[language=Python, caption={Configuración de la operación de composición según herramienta activa}, label={lst:composicion}]
if (esBorrador) {
  contextRef.current.globalCompositeOperation = 'destination-out';
  contextRef.current.lineWidth = grosurGoma;
} else {
  contextRef.current.globalCompositeOperation = 'source-over';
  contextRef.current.strokeStyle = 'black';
  contextRef.current.lineWidth   = grosurLapiz;
}
\end{lstlisting}

Al finalizar cada trazo, la operación de composición se restaura a
\texttt{'source-over'} y se guarda el estado actual del canvas.

\subsubsection{Historial de estados, pantalla completa y envío}

\textbf{Deshacer.} Cada vez que el alumno levanta el lápiz o el dedo,
el contenido completo del canvas se serializa como \textit{data URL}
y se agrega a un arreglo de estados. La acción de deshacer extrae el
último elemento de ese arreglo y restaura el estado anterior
redibujando la imagen serializada sobre el canvas; una condición de
guarda impide retroceder más allá del estado inicial (lienzo en blanco).

\textbf{Modo de pantalla completa.} El lienzo puede expandirse a las
dimensiones de la ventana para ofrecer mayor área de escritura. Dado
que redimensionar el canvas borra su contenido, el sistema verifica
si existen trazos en el historial antes de aplicar el cambio; de ser
así, muestra un diálogo de confirmación que advierte al alumno sobre
la pérdida del contenido. Una vez confirmado, un \texttt{useEffect}
dependiente del estado de pantalla completa reinicializa el canvas con
un retardo de 150\,ms para que React aplique los cambios de CSS antes
de releer las dimensiones del contenedor.

\textbf{Exportación y envío.} Al concluir, el lienzo se convierte a 
PNG mediante \texttt{canvas.toDataURL(}\allowbreak\texttt{'image/png')} 
y se envía al mismo endpoint de registro de intentos, siguiendo el mismo 
esquema de autenticación por token JWT que \textit{CapturaEscritura}. 
A diferencia de este último, que presenta los errores de forma persistente, 
\textit{LienzoDigital} utiliza notificaciones temporales que desaparecen 
automáticamente a los cuatro segundos para no obstruir el área de trabajo 
de forma prolongada.

% ------------------------------------------------------------
\subsection{Contrato compartido con el \textit{backend}}
% ------------------------------------------------------------

Ambos componentes del módulo de captura convergen en el mismo
contrato de comunicación hacia el servidor. El cuerpo de la petición
es un objeto JSON con dos campos: \texttt{id\_ejercicio\_tutor},
que identifica la asignación a la que corresponde el intento, e
\texttt{imagen\_codificada}, que contiene la imagen serializada como
\textit{data URL}. La cabecera \texttt{Authorization} transporta el
token JWT para que el middleware del servidor pueda identificar al
alumno y verificar que el ejercicio le fue asignado antes de persistir
el intento:

\begin{lstlisting}[language=JSON, caption={Estructura del cuerpo de la petición POST al endpoint de registro de intentos}, label={lst:payload_intento}]
{
  "id_ejercicio_tutor": 42,
  "imagen_codificada":
    "data:image/jpeg;base64,/9j/4AAQSkZJRgAB..."
}
\end{lstlisting}

Esta uniformidad del contrato permite que el backend procese los
intentos sin distinguir el origen de la imagen (fotografía o lienzo
digital), simplificando la lógica del endpoint de registro y
favoreciendo la extensibilidad del módulo ante futuros modos de
captura adicionales.